\documentclass{standalone}
\usepackage{circuitikz}
\usetikzlibrary{calc}
\definecolor{circuitnotemark}{HTML}{FE00FE}
\begin{document}
\begin{circuitikz}[american, line width=0.8pt]
\ctikzset{bipoles/length=1.2cm}
\foreach \x in {0,2,4,6,8} {\foreach \y in {0,-2,-4,-6,-8,-10} {\fill[gray, opacity=0.35] (\x,\y) circle (1.1pt);}}
\node[gray, font=\scriptsize] at (0,0.84) {1};
\node[gray, font=\scriptsize] at (2,0.84) {2};
\node[gray, font=\scriptsize] at (4,0.84) {3};
\node[gray, font=\scriptsize] at (6,0.84) {4};
\node[gray, font=\scriptsize] at (8,0.84) {5};
\node[gray, font=\scriptsize, anchor=east] at (-0.84,0) {a};
\node[gray, font=\scriptsize, anchor=east] at (-0.84,-2) {b};
\node[gray, font=\scriptsize, anchor=east] at (-0.84,-4) {c};
\node[gray, font=\scriptsize, anchor=east] at (-0.84,-6) {d};
\node[gray, font=\scriptsize, anchor=east] at (-0.84,-8) {e};
\node[gray, font=\scriptsize, anchor=east] at (-0.84,-10) {f};
\coordinate (b1) at (0,-2);
\coordinate (b5) at (8,-2);
\coordinate (d1) at (0,-6);
\coordinate (d5) at (8,-6);
\coordinate (f1) at (0,-10);
\coordinate (f5) at (8,-10);
\coordinate (a3) at (4,0);
\coordinate (c3) at (4,-4);
\coordinate (e3) at (4,-8);
\draw (b1) to[potentiometer, n=part-P1, l_=$P_{1}$, a^=$10\,\mathrm{k}\Omega$] (b5); % line 2
\draw (d1) to[thyristor, n=part-T1, l_=$T_{1}$] (d5); % line 3
\draw (f1) to[triac, n=part-T2, l_=$T_{2}$] (f5); % line 4
\draw (part-P1.wiper) -- (a3); % line 6
\draw (part-T1.gate) |- (c3); % line 7
\draw (part-T2.gate) |- (e3); % line 8
\path (12,-0.226) rectangle (16.568,0.226); % line 10
\node[anchor=west, circuitnotemark, font=\footnotesize] at (12,0) {X}; % line 10
\path (12,-2.226) rectangle (15.215,-1.774); % line 11
\node[anchor=west, circuitnotemark, font=\footnotesize] at (12,-2) {X}; % line 11
\path (12,-4.226) rectangle (14.538,-3.774); % line 12
\node[anchor=west, circuitnotemark, font=\footnotesize] at (12,-4) {X}; % line 12
\end{circuitikz}
\end{document}
